
\subsection{The Authority of the Synod/Church Body and the Free Congregation.}

Source: Editorial articles in “Folketbladet” for January 6, 1897; April 20 and 27; and May 25, 1898. — Ed.


\begin{quote}
\emph{Note: This section is currently represented by an editorial summary. 
A full translation of the original text is planned and contributions are welcome.}
\end{quote}

This text is a sustained critique of ecclesiastical centralization and majority rule within church organizations, especially as practiced in Den forenede Kirke in the late 19th century. The author argues that church “societies” (samfund) have come to be treated as quasi-divine authorities that demand obedience from local congregations, a mindset reflected in phrases like “under the synod” or “under the society.” This, the author contends, mirrors a papist (Roman Catholic–style) understanding of authority, where decisions made by assemblies are regarded as direct expressions of God’s will, and dissent is treated as sin.

According to this logic, annual church meetings (årsmøder) function like councils or episcopal authorities: their majority decisions are equated with God’s will, and those who disagree are seen as opposing divine order. Over time, this has ingrained the idea that the church organization itself is a divine authority on earth, enforced not only by rhetoric but also by exclusion and discipline. The author criticizes this as a dangerous substitution of human authority for God’s Word.

The text contrasts this view with a biblical understanding of the church. In Scripture, the local congregation (menigheden) is the highest form of Christian community. It answers directly to God and His Word, not to synods, constitutions, or church officials. No church body has the right to impose obligations or release people from obligations apart from Scripture. Cooperation among congregations is possible and desirable, but only when it occurs freely, in love, and without coercion.

A major portion of the text focuses on what the author calls the “misuse of majority power.” While acknowledging that congregations may freely enter cooperative agreements or constitutions, the author insists that such arrangements are not divinely mandated and must never override conscience or Scripture. The constitution of Den forenede Kirke grants extensive power to annual meetings, including the ability to decide matters by a simple majority of those present, with no protection for minorities. This, the author argues, is not inherently wrong but is extremely risky—and in practice has been gravely abused.

Concrete examples are given: the 1893 splitting of Augsburg Seminary, forced removals of professors and students, the division of congregations in Minneapolis, and the expulsion of minority-aligned congregations and delegates in 1895–1896. These actions illustrate how the majority used its power not merely to govern common affairs but to silence, exclude, and dispossess dissenters. By expelling minorities, the church denied them even the basic right to continue working peacefully for their convictions—a right normally preserved even in secular political systems.

The problem is intensified by the church’s dual nature as both a religious body and a legal corporation holding property. By expelling congregations, the majority could also deprive them of property, something even state churches could not do. This fusion of spiritual authority and corporate power allows “right to become wrong, and wrong to become right.”

The author concludes that the abuses of majority rule are so severe that a new starting point is necessary: congregational freedom must come first, even if institutional power diminishes. Finally, the text argues that free congregations can cooperate. True unity arises not from enforced structures but from the shared work of God’s Spirit. When congregations are spiritually alive, independent, and responsible, cooperation—such as in mission work—will naturally follow. Freedom, knowledge, and spiritual maturity do not hinder cooperation; they are its true foundation.